% Generated draft: requires mathematical and source review.
\documentclass[11pt]{article}
\usepackage[margin=25mm]{geometry}
\usepackage{fontspec}
\setmainfont{Latin Modern Roman}
\usepackage{amsmath,amssymb,mathtools,unicode-math}
\setmathfont{Latin Modern Math}
\usepackage{enumitem,needspace,xurl,hyperref}
\setlength{\parindent}{0pt}
\setlength{\parskip}{5pt}
\setlength{\emergencystretch}{3em}
\newcommand{\N}{\mathbb{N}}
\newcommand{\Z}{\mathbb{Z}}
\newcommand{\Q}{\mathbb{Q}}
\newcommand{\R}{\mathbb{R}}
\newcommand{\C}{\mathbb{C}}
\newcommand{\F}{\mathbb{F}}
\newcommand{\A}{\mathbb{A}}
\newcommand{\PP}{\mathbb{P}}
\newcommand{\E}{\mathbb{E}}
\newcommand{\eps}{\varepsilon}
\newcommand{\dd}{\,\mathrm{d}}
\begin{document}
\section*{UnsolvedMath ID: 10400112}
\textbf{Source problem code:} catalogue entry
\subsection{Definitions and mathematical statement}
The BMW algebra is the Birman--Murakami--Wenzl braid/tangle algebra, with parameters and coefficient ring that must be specified. It has braid generators subject to braid and Kauffman-type skein/tangle relations. Write \(\operatorname{BMW}_n\) for the source's version on \(n\) strands; the supplied extract does not give its presentation or parameter regime. Over a field \(F\), a nonzero representation \(\rho:\operatorname{BMW}_n\to\operatorname{End}_F(V)\) is irreducible when \(V\) has no nonzero proper invariant subspace.

S. J. Bigelow's Problem 6.4 asks to generalize Lawrence's construction so as to obtain the irreducible representations of the BMW algebra. Lawrence's construction refers to geometric braid-group representations built from homology of configuration spaces associated with a punctured disk, using covering or local-coefficient data. The exact construction to be generalized is not included in the record.

A precise version of the requested goal would specify geometric homology data and a braid action, prove that the action satisfies the defining BMW relations, and identify the resulting BMW modules with the intended irreducible representations. The source asks for the irreducible representations collectively, not merely one faithful braid representation. However, it does not state the strand range, coefficient field, generic versus specialized parameters, semisimplicity assumptions, or which homology/subquotient constructions are admissible. These choices affect what ``the irreducible representations'' means and must be recovered rather than silently restricted to a convenient regime.
\subsection{Short English statement}
Extend Lawrence's geometric braid-representation construction to realize the irreducible BMW-algebra representations. The source leaves the algebra's parameters and the precise homological construction unspecified. Its supplied open-triage summary reports no clear resolution.
\subsection{Sources}
\begin{itemize}
\item Local UnsolvedMath registry, permanent ID \(10400112\), exact source code \texttt{catalogue entry}, Problem 6.4, attributed to S. J. Bigelow. The supplied title is truncated after ``(S.J.''
\item Ohtsuki, \emph{Problems on invariants of knots and 3-manifolds} (2002), Problem 6.4, PDF page 97: \url{https://msp.org/gtm/2002/04/gtm-2002-04-024s.pdf}. Bold printed headings distinguish 231 numbered items from references in surrounding exposition.
\item Extraction pdf-font-structure, accessed 2026-07-29. Source evidence NEEDS\_REVIEW presents an open problem and calls for a current primary-literature audit. Supplied research classification OPEN-TRIAGE calls this Lawrence-to-BMW generalization not obviously resolved. No separate reference for Lawrence's construction or the BMW presentation is supplied.
\item Dataset attribution: Ulam AI / UnsolvedMath Contributors, CC BY 4.0. Source rights NEEDS\_REVIEW: publicly accessible author/publisher copy, redistribution terms requiring review.
\end{itemize}
\end{document}
